\chapter{Field Operations Guide}
\label{ch:field-operations}

This chapter provides essential guidelines and operational procedures for using the Mountain Climber Health \& GPS Tracker system in real-world mountaineering environments. Proper preparation and adherence to these protocols are critical for ensuring reliable communication, accurate health monitoring, and system longevity during an expedition.

\section{Pre-Climb Preparation Checklist}
\label{sec:pre-climb-checklist}

Before commencing any expedition, ensure that all system components are fully operational. A systematic pre-climb check minimizes the risk of equipment failure in critical situations.

\begin{itemize}
    \item \textbf{Charge All Devices:} Verify that the Climber Main Device, Wristband, Repeater nodes, and Base Camp Unit are charged to 100\%. Check the battery voltage via the mobile application or dashboard.
    \item \textbf{Test BLE Pairing:} Pair the Climber Main Device with the Wristband and mobile application. Confirm that real-time heart rate and SpO2 data are successfully transmitted and displayed on the app.
    \item \textbf{Verify GPS Fix:} Take the Climber Main Device outdoors to establish a clear line of sight to the sky. Confirm that the NEO-6M GPS module acquires a valid location fix within 5 minutes (indicated by the double beep and dashboard update).
    \item \textbf{Test LoRa Communication Range:} Conduct a brief range test between the Climber Main Device and the Base Camp Unit (or through a Repeater) to ensure data packets are successfully received. Verify that the Received Signal Strength Indicator (RSSI) is within acceptable limits.
    \item \textbf{Pack Spares and Tools:} Ensure that spare 18650 Li-ion batteries, a portable power bank, micro-USB cables, and necessary tools for battery replacement are included in your gear.
\end{itemize}

\section{Deployment Strategy for Repeaters}
\label{sec:repeater-deployment}

In mountainous terrain, line-of-sight (LOS) communication is frequently obstructed by ridges, dense forests, and deep valleys. Strategic placement of Repeater nodes is essential to maintain network connectivity.

\begin{itemize}
    \item \textbf{Line-of-Sight Placement:} Position repeaters at elevated, unobstructed locations. The signal propagates best when both the transmitting and receiving antennas have a clear visual path.
    \item \textbf{Elevation Considerations:} Deploy repeaters on ridges or high points rather than in valleys. Avoid placing them directly on the ground; elevate them on poles or attach them to robust structures or trees.
    \item \textbf{Antenna Orientation:} Ensure the LoRa antennas are oriented vertically to maximize the omnidirectional radiation pattern.
    \item \textbf{Weather Protection:} Secure the repeater enclosures against harsh weather conditions, ensuring that water cannot penetrate the casing or antenna connectors.
\end{itemize}

\section{Communication Range Tips}
\label{sec:comm-range-tips}

The SX1278 LoRa transceiver operates at 433MHz with Spreading Factor 7 (SF7) and 125kHz Bandwidth. While it is designed for long-range communication, environmental factors significantly impact performance.

\begin{tipbox}
To maximize LoRa range, keep the Climber Main Device attached to the upper part of the backpack, ensuring the antenna is exposed and points skyward.
\end{tipbox}

\begin{itemize}
    \item \textbf{Clear Line of Sight (LOS):} Under ideal LOS conditions, the communication range can reach 2 to 5 kilometers.
    \item \textbf{Obstructed Environments:} In dense forests, deep valleys, or rocky canyons, the range may be reduced to 500 meters or less. Utilize repeaters to bridge these gaps.
    \item \textbf{Body Interference:} The human body absorbs RF signals. Avoid covering the main device's antenna with your body, metallic objects, or wet clothing.
\end{itemize}

\section{Battery Management on Multi-Day Climbs}
\label{sec:battery-management}

Managing battery life is crucial for multi-day expeditions. The system relies on 18650 Li-ion batteries, and power consumption must be optimized.

\begin{itemize}
    \item \textbf{Update Intervals:} If battery levels drop below 30\%, adjust the transmission frequency via the mobile app to send GPS and health updates less frequently (e.g., every 5 minutes instead of every 1 minute).
    \item \textbf{Cold Weather Performance:} Lithium-ion batteries degrade in extremely cold temperatures. Keep the main device reasonably insulated.
    \item \textbf{Power Down When Stationary:} If resting for extended periods at a safe camp, consider powering down the system or switching to a low-power standby mode to conserve battery.
\end{itemize}

\section{Emergency SOS Procedure}
\label{sec:emergency-sos}

The SOS functionality is the most critical feature of the system, designed to alert the Base Camp immediately in case of an emergency.

\begin{cautionbox}
Do not trigger the SOS alert unless there is a genuine emergency requiring immediate assistance. False alarms can lead to unnecessary deployment of rescue resources.
\end{cautionbox}

\textbf{How to Trigger SOS:}
\begin{enumerate}
    \item Locate the dedicated physical SOS button on the Climber Main Device or use the digital SOS button in the mobile application.
    \item Press and hold the physical button for 3 seconds, or slide the confirmation bar in the mobile app.
\end{enumerate}

\textbf{What Happens When Triggered:}
\begin{itemize}
    \item The Climber Main Device immediately overrides standard transmission schedules and broadcasts an emergency LoRa packet with the highest priority.
    \item The packet contains the precise GPS coordinates and the latest health vitals (Heart Rate and SpO2).
    \item The local buzzer on the Climber Main Device will emit a continuous emergency tone, and the LED will flash red rapidly to confirm activation and aid visual location.
    \item The Base Camp dashboard will trigger severe visual and audible alarms, displaying the climber's exact location on the map.
\end{itemize}
